\subsection{线性表示的正则元}

引理 1. 设 M 是 k 上的一个解析流形，\(a = (a_{0}, \ldots, a_{n-1}, a_{n} = 1)\) 是 M 上的一个解析函数序列。对所有 \(x \in M\)，设 \(r_{a}(x)\) 是那些 \(i \in [0, n]\)（它们满足 \(a_{j}(x) = 0\)，其中 j \textless{} i）的上确界，并设 \(r_{a}^{0}(x)\) 是那些 \(i \in [0, n]\)（对它们，\(a_{j}\) 在 x 的一个邻域上为零，其中 j \textless{} i）的上确界。

\begin{enumerate}
\def\labelenumi{(\roman{enumi})}
\item
  函数 \(r_a\) 是上半连续的。
\item
  对所有 \(x \in \mathrm{M}\)，\(r_a^0(x) = \liminf_{y \to x} r_a(y)\)。
\item
  函数 \(r_a^0\) 是局部常值的。
\item
  点 \(x \in M\)（满足 \(r_{a}^{0}(x) = r_{a}(x)\)）所成的集，就是在 \(r_{a}\) 于其一个邻域上取常值的 M 的点所成的集。这是 M 的一个稠密开子集。若 k = C 且 M 有限维并连通，则它还是开且连通的。
\item
  若 \(r_a(x) = i\)，则 \(a_i(x) \neq 0\)，且对所有 \(y\)（\(x\) 的一个邻域中的点），有 \(a_i(y) \neq 0\)，从而 \(r_a(y) \leq i\)。
\item
  若 \(r_{a}^{0}(x) = i\)，则函数 \(a_{0}, \ldots, a_{i-1}\) 在 x 的一个邻域上为零，且对该邻域中的任何 y，\(r_{a}(y) \geq i\)。因而 \(\liminf_{y \to x} r_{a}(y) \geq i\)。x 的每个邻域都含一点 y，使 \(a_{i}(y) \neq 0\)，从而 \(r_{a}(y) \leq i\)。于是 \(\liminf_{y \to x} r_{a}(y) = i\)。
\item
  设 \(i = r_a^0 (x)\)，并设 \(V\) 是 \(x\) 的一个邻域，使得 \(a_{j}(y) = 0\) 对所有 \(y\in V\) 和所有 \(j < i\) 成立。则 \(x\in M - Z\)，其中 \(Z\) 表示 \(M\) 中在其一个邻域上函数 \(a_{i}\) 为零的点所成的集。由于 \(Z\) 在 \(M\) 中是闭的（《微分与解析流形·结果》，5.3.5），\(V\cap (M - Z)\) 是 \(x\) 的一个邻域。对此邻域中的每一点 \(y\)，有 \(r_a^0 (y) = i\)。
\item
  函数 \(r_{a}-r_{a}^{0}\) 是上半连续的，且它在任一点的值都 \(\geq0\)。若 \(r_{a}(x)=r_{a}^{0}(x)\)，则 \(r_{a}-r_{a}^{0}\) 在 x 的一个邻域上为零，由 (iii) 这表明 \(r_{a}\) 在 x 的一个邻域上为常数。反之，若 \(r_{a}\) 在 x 的一个邻域上为常数，则由 (ii) 有 \(r_{a}^{0}(x)=r_{a}(x)\)。因此，点 \(x\in M\)（满足 \(r_{a}^{0}(x)=r_{a}(x)\)）所成的集是 M 的一个开子集 \(\Omega\)。若 \(x\in M\) 且 \(r_{a}^{0}(x)<r_{a}(x)\)，则 x 的每个邻域都含一点 y，使得 \(r_{a}(y)<r_{a}(x)\) 且 \(r_{a}^{0}(y)=r_{a}^{0}(x)\)。于是 x 的每个邻域都含一点 y，使得
\end{enumerate}

\[
r _ {a} (y) - r _ {a} ^ {0} (y) <   r _ {a} (x) - r _ {a} ^ {0} (x).
\]

由此推出 \(\Omega\) 在 M 中稠密。

若 M 连通且 p 是 \(r_{a}^{0}\) 在 M 上的值，则 \(\Omega\) 的点就是那些点 \(x \in M\)：它们满足 \(a_{p}(x) \neq 0\)。若 k = C，则由附录 II 的引理 3 推出 \(\Omega\) 连通。

设 \(\rho\) 是 \(G\) 在一个向量空间 \(V\)（维数 \(n\) 有限、在 \(k\) 上）上的解析线性表示。置

\[
\det (\mathrm{T} - \rho (g) + 1) = a _ {0} (g) + a _ {1} (g) \mathrm{T} + \dots + a _ {n - 1} (g) \mathrm{T} ^ {n - 1} + \mathrm{T} ^ {n}.
\]

函数 \(r_a\) 与 \(r_a^0\)（与序列 \((a_0, a_1, \ldots, a_{n-1}, 1)\) 相伴者）将分别记作 \(r_\rho\) 与 \(r_\rho^0\)。于是，对所有 \(g \in G\)，

\[
\begin{array}{l} r _ {\rho} (g) = \dim \mathrm{V} ^ {1} (\rho (g)) \\ r _ {\rho} ^ {0} (g) = \liminf _ {g ^ {\prime} \to g} \dim \mathrm{V} ^ {1} (\rho (g ^ {\prime})). \end{array}
\]

引理 2. 设 \(0 \to V' \to V \to V'' \to 0\) 是一个 G-模的正合列，分别由 G 的解析线性表示 \(\rho', \rho, \rho''\) 所定义。则：

\[
r _ {\rho} = r _ {\rho^ {\prime}} + r _ {\rho^ {\prime \prime}}, \text { and } r _ {\rho} ^ {0} = r _ {\rho^ {\prime}} ^ {0} + r _ {\rho^ {\prime \prime}} ^ {0}.
\]

实际上，对所有 \(g \in \mathbf{G}\)，存在一个正合列（§1，第 1 节，定理 1 的推论 3）

\[
0 \to (\mathrm{V} ^ {\prime}) ^ {1} (\rho^ {\prime} (g)) \to \mathrm{V} ^ {1} (\rho (g)) \to (\mathrm{V} ^ {\prime \prime}) ^ {1} (\rho^ {\prime \prime} (g)) \to 0,
\]

这就证明了第一个断言。第二个断言由此推出，因为由引理 1 (iv)，在 G 的一个稠密开子集上 \(r_{\rho}^{0} = r_{\rho}, r_{\rho'}^{0} = r_{\rho'}\) 且 \(r_{\rho''}^{0} = r_{\rho''}\)。

定义 1. 元素 \(g \in \mathbf{G}\) 称为关于线性表示 \(\rho\) 的正则元，如果 \(r_{\rho}(g) = r_{\rho}^{0}(g)\)。

命题 1. 对于 G 的解析线性表示 \(\rho\)，其正则点就是 G 中在其一个邻域上 \(r_{\rho}\) 为常数的那些点。它们构成 G 的一个稠密开子集。若 k = C 且 G 连通，则 \(\rho\) 的正则点所成的集是连通的。

这由引理 1 (iv) 得出。

注记. 设 \(G^{*}\) 是 G 的一个开子群。元素 \(a \in G^{*}\) 是 G 的关于 G 的线性表示 \(\rho\) 的正则元，当且仅当它是 \(G^{*}\) 的关于线性表示 \(\rho|G^{*}\) 的正则元。

